\section{Benchmarks and Setup}
当前稿件的默认故事线统一采用 \textbf{长时程} 口径，不再只覆盖 onset 附近一小段时间。每个 benchmark 至少跨越 onset、peak compression、release 以及 release 后的一段自由段，使得 peak force、impulse、energy drift 与 rebound velocity 等指标具有真正的解释力。

正文主线 benchmark 包含三组：
\begin{itemize}
    \item \textbf{Analytic flat}：面积主导的解析对照组；
    \item \textbf{Analytic sphere}：曲率主导的解析对照组；
    \item \textbf{Native-band flat}：SDF-native 主链的真实动态 benchmark。
\end{itemize}
此外，本文再补入一组 \textbf{6-DoF complex rigid-body benchmark}，用于验证方法在非对称复合几何、平动--转动耦合、分布式接触 wrench 以及 active-set 拓扑传递下的外推能力。

统一报告的主要指标包括：peak force error、impulse error、energy drift、release timing error、max penetration error 与 rebound velocity error。对于 native-band benchmark，还同时记录 mean candidate cells、mean recompute cells、predictor/corrector Jaccard 等效率与连续性指标；对于 6-DoF benchmark，则进一步报告位置/姿态/力/力矩 RMS 误差、peak wrench 误差以及 active-component-count 误差。
